\documentclass[11pt]{article}
\usepackage{fontspec}
\usepackage[french,english]{babel}
\usepackage{amsmath}
\usepackage{amssymb}
\usepackage{amsthm}
\usepackage{mathrsfs}
\usepackage{tikz-cd}
\usepackage{enumitem}
\usetikzlibrary{babel,decorations.pathreplacing}

\newcommand{\printedpage}[1]{\marginpar{\footnotesize p.~#1}}
\newcommand{\msp}[1]{\marginpar{\footnotesize [#1]}}
\newcommand{\M}{\mathcal{M}}
\newcommand{\Mc}{\mathcal{M}_{\mathrm{c}}}

\begin{document}

\printedpage{19}
% head: A. Grothendieck : Dérivateurs, Ch. XVII, Édition des Universités de Montpellier II et Paris VII

\[
\begin{tikzcd}
X_{I_\alpha} \arrow[r] \arrow[d] & X_{I_{\alpha_0}} \arrow[d] \\
Y_{I_\alpha} \arrow[r] & Y_{I_\alpha,I_{\alpha_0}} \rlap{${}= Z_\alpha$} \arrow[d] \\
& Y_{I_{\alpha_0}} \rlap{\ ,}
\end{tikzcd}
\]
we have $Z_0 = Y_{I_\alpha,I_{\alpha_0}}$, $Z_{\alpha_0} = Y_{I_\alpha,I_{\alpha_0}}$, $Z_\alpha = \varinjlim\limits_{\beta<\alpha} Z_\beta$ if $\alpha$ \msp{34} a limit ordinal. So it remains to determine the $v_{\alpha,\alpha+1} : Z_\alpha \longrightarrow Z_{\alpha+1}$ (for $\alpha < \alpha_0$), which we see from the diagram
\[
\begin{tikzcd}
X_{I_\alpha} \arrow[r] \arrow[d] & X_{I_{\alpha+1}} \arrow[r] \arrow[d] & X_{I_{\alpha_0}} \rlap{${}= Z_0$} \arrow[d] \\
Y_{I_\alpha} \arrow[r] & Y_{I_\alpha,I_{\alpha+1}} \arrow[r] \arrow[d, "u_{I_\alpha,I_{\alpha+1}}"] & Y_{I_\alpha,I_{\alpha_0}} \rlap{${}= Z_\alpha$} \arrow[d, "v_{\alpha,\alpha+1}"] \\
& Y_{I_{\alpha+1}} \arrow[r] & Y_{I_{\alpha+1},I_{\alpha_0}} \rlap{${}= Z_{\alpha+1}\ .$}
\end{tikzcd}
\]

Let us resume the proof of 2.1. We are going to prove 1° by the two preceding lemmas. Let us note that when the space underlying $I$ is noetherian, then the transfinite construction is necessarily finite, so we do not need stability of $C$ under infinite transfinite composites.

\msp{35}

Let us prove 2°. We therefore assume that for every $i_0 \in I$, we have $u_{I_{<\alpha_0},I/i_0} \in C$, we want to deduce from this that $u_\bullet \in C_I$, i.e. $u_{J,I/i_0} \in C$ for \textit{every} open set.

\noindent\textbf{Lemma 2.8.} (Under the artinian hypothesis on $I$.) \textit{Let $P \subset I$ be the set of $i_0 \in I$ such that for $u_\bullet|I/i_0$ to be in $C_{I/i_0}$, i.e. that for every $i \leq i_0$, and every open set $K$ of $I/i$, we have $u_{K,I/i} \in C$. Let $i_0 \in I$ be such that $I_{<i_0} \subset P$. Then $i_0 \in P$.}

\noindent\textsc{Proof}. It remains to prove that if $K \subset I/i_0 = J$ is an open set, then $u_{K,J} \in C$. We may assume $K \neq I/i_0$, i.e. $K \subset I_{<i_0} \overset{\text{def}}{=} J'$. But then by the initial hypothesis on $u_\bullet$, we have $u_{J',J} \in C$, and on the other hand, by the hypothesis on $i_0$, $u_\bullet|J'$ is in $C_{J'}$. So by theorem 2.1. 1°, we have $u_{K,J'} \in C$. So by transitivity $u_{K,J} \in C$, q.e.d.

\noindent\textbf{Corollary 2.9.} \textit{If $I$ is artinian, then $P = I$, i.e. $u_\bullet \in C_I$.}

This therefore proves 2.1. 2°. Let us note that $\mathrm{Esp}(I)$ is noetherian, the same is true of every open set (such as $J'$ above), so \msp{36} for the conclusion of the lemma, no need to assume $C$ stable under infinite composition.

\printedpage{20}
% head: A. Grothendieck : Dérivateurs, Ch. XVII, Édition des Universités de Montpellier II et Paris VII

Let us prove 2.1. 3°. We have a) $\Longrightarrow$ b), for by 2.1. 1°, $u_\bullet \in (TC)_I$ implies that for $i \in I$, $u_i : X_i \longrightarrow Y_i$ is in $TC$ (cf. cor. 2.1 a). So it remains to prove b) $\Longrightarrow$ a), i.e. if $u_\bullet \in C \cap W_{I^\circ}$, then $u_\bullet \in (TC)_I$, i.e. the $u_{K,I/i_0}$ are in $TC$. We already know that $u_{K,I/i_0} \in C$ by hypothesis $u_\bullet \in C_I$, so we must prove that it is in $W$. We are reduced to proving that $u_K \in W$, i.e. $u_K \in TC$, for by the diagram (where $J = I/i_0$)
\[
\begin{tikzcd}
X_K \arrow[r] \arrow[d, "u_K"'] \arrow[dr, phantom, "\text{cocart.}"] & X_J \arrow[d, "\tilde{u}_{K,J}"'] \arrow[dd, bend left=50, "u_J"] \\
Y_K \arrow[r] & Y_{K,J} \arrow[d, "u_{K,J}"'] \\
& Y_J \rlap{\ ,}
\end{tikzcd}
\]
$u_K \in TC$ implies $\tilde{u}_{K,J} \in TC$, and since $u_J \simeq u_{i_0} \in W$ by hypothesis, it indeed follows that $u_{K,J} \in W$ by saturation. But we are going to go about it differently:

\msp{37}

\noindent\textbf{Lemma 2.9.} \textit{$I$ an arbitrary ordered set, $C$, $W$ as in th.} 2.1, 1° \textit{and} 3°, \textit{$u_\bullet : X_\bullet \longrightarrow Y_\bullet$ in $C_I \cap W_{I^\circ}$. Let $P$ be the set of $i_0 \in I$ such that $u_\bullet|I/i_0$ is in $(TC)_{I/i_0}$. If $i_0 \in I$ is such that $I_{<i_0} \subset P$, then $i_0 \in P$.}

Indeed, it remains to prove that for every open set $K$ of $J = I/i_0$, $K \neq J$, we have $u_{K,J} \in TC$. But for this, it suffices to see that $u_K, u_J \in W$ (as we had just seen). Now $u_J \simeq u_{i_0} \in W$ by hypothesis. On the other hand $u_\bullet|K \in (TC_K)$ since $K \subset I_{<i_0} \subset P$, so by theorem 2.1, 1°, we conclude that $u_K \in TC$.

If now we assume $I$ artinian, it follows from 3.9 that $P = I$, so $u_\bullet \in (TC)_I$, q.e.d. % [sic] 3.9 -> 2.9, the lemma just stated is Lemme 2.9

This completes the proof of theorem 2.1, and also of corollary 2.5, seeing that for the proof of lemma 2.9 as well, the infinite hypotheses on $C$, $C \cap W$ are unnecessary if $\mathrm{Top}(I)$ is noetherian, hence its open sets are also noetherian.

\noindent\textbf{Remarks.}
\begin{enumerate}[leftmargin=2.5em]
\item[1)] In 2.1. 2°, the hypothesis $I$ artinian is not superfluous, as we see by taking $I = \mathbf{N}^{\mathrm{opp}} \simeq \mathbf{Z}^{\leq 0}$, \msp{38} and for an arrow % [sic] n'est par superflue -> n'est pas superflue, negation "pas" intended
\[ \emptyset_{\M(I)} \longrightarrow Y_\bullet\ . \]

To say that this arrow is in $C_I$ means that the transition morphisms $Y_i \xrightarrow{u_i} Y_{i+1}$ are in $C$, \textit{and} that the $Y_i$ are $C$-cofibrant, i.e. $\emptyset_{\M} \longrightarrow Y_i$ is in $C$, for every $i$. But the condition considered in 2.1. 1° is only that the $u_i$ are in $C$. So by taking for $Y_\bullet$ a constant direct system with value $Y$, with $Y$ not $C$-cofibrant, i.e. $\emptyset_{\M} \longrightarrow Y$ not in $C$, we have a counterexample.

With the same $I$, we see that $u_\bullet : X_\bullet \longrightarrow Y_\bullet$ is in $C$ if and only if
\begin{enumerate}[leftmargin=2.5em]
\item[a)] the $u_i \in C$, and
\item[b)] the $u_{i,i+1} : Y_{i,i+1} \longrightarrow Y_{i+1}$ are in $C$.
\end{enumerate}

\printedpage{21}
% head: A. Grothendieck : Dérivateurs, Ch. XVII, Édition des Universités de Montpellier II et Paris VII

If moreover the $u_i$ are in $W$, hence in $TC$, we see as above that $u_{i,i+1} \in TC$ for every $i$, hence $u_\bullet \in (TC)_I$. So here we still have
\[ (TC)_I \ = \ C_I \cap W_{I^\circ}\ , \]
although $I$ is not artinian. I do not know whether the artinian hypothesis in 2.1. 3° is superfluous or not.
\item[2)] Let $u_\bullet \in \mathrm{Fl}(\M(I^\circ))$\footnote[13]{$I^\circ = I^{\mathrm{opp}}$.}, so $u_\bullet$ is an object of $\underline{\mathrm{Fl}}(\M(I^\circ)) \simeq \M(I^\circ \times \boldsymbol{\Delta}^1) = \M((I \times (\boldsymbol{\Delta}^1)^{\mathrm{opp}})^{\mathrm{opp}})$. We check that for this object to be $C$-cofibrant, it is necessary and sufficient that
\begin{enumerate}[leftmargin=2.5em]
\item[1°)] $X_\bullet$ be $C$-cofibrant, and
\item[2°)] $u_\bullet \in C_I$ (which implies that $Y_\bullet$ is also cofibrant).
\end{enumerate}
\end{enumerate}
\msp{39}

\section*{3\quad Categories of models with cofibrations}

Let $\M$ be a category equipped with two sets of arrows
\[ W,\ C\ \subset\ \mathrm{Fl}(\M)\ . \]
We assume that $\M$ contains an initial object, and:
\begin{description}[leftmargin=3em, labelwidth=2.5em]
\item[Cof$_0$] $W$ is mildly saturated, i.e. contains the isomorphisms and if $f$, $g$ are composable, and if two among $f$, $g$, $gf$ are in $W$, the third too.
\item[Cof$_1$] $C$ contains the isomorphisms, is stable under composition and under cobase change.
\item[Cof$_2$] $TC \overset{\text{def}}{=} C \cap W$ is stable under cobase change.
\item[Cof$_3$] Every arrow $f$ in $\M$ is written
\[ f \ = \ pi \]
with $i \in C$, $p \in W$.
\end{description}
These are the axioms of K. S. \textsc{Brown} (dual form), who moreover assumes that the objects of $\M$ are $C$-cofibrant, i.e. $\emptyset_{\M} \longrightarrow X$ is in $C$. I will not make this assumption. The subcategory $\Mc$ of $\M$ formed of the \msp{40} $C$-cofibrant objects satisfies the axioms of \textsc{Brown}, moreover one proves that
\[ \mathrm{Ho}_W(\Mc) \xrightarrow{\ \sim\ } \mathrm{Ho}_W(\M) \]
is an equivalence. Thus the results of \textsc{Brown} concerning $\Mc$, equipped with the sets of arrows induced by $W$, $C$, are applicable to $\mathrm{Ho}_W(\M)$ and to its avatars (such as $\mathrm{Ho}_W(S\backslash\M)$, $\mathrm{Ho}_W(\M/T)$ and $\mathrm{Ho}_W(S\backslash\underset{u_0}{\M}/T)$, \ where $u_0 : S \longrightarrow T$ is given). It turns out that for the most part, the standard homotopy results (in particular exact sequence of cofibrations relating to the homotopy operations ``on the right'' (or should one say on the left? yet these are questions of right exactness, not left exactness)) are valid. I will have to return to this in greater detail later on. In particular, there is a good theory of $W$-cocartesian squares (and \textit{in} $\Mc$\textit{,} there are enough $W$-cofibrations\footnote[14]{More precisely, by K. S. \textsc{Brown}, \textit{in} $\Mc$\textit{,} the objects of $C$ are $W_{\mathrm{c}}$-cofibrations.}, \ldots),
\printedpage{22}
% head: A. Grothendieck : Dérivateurs, Ch. XVII, Édition des Universités de Montpellier II et Paris VII
\msp{41}
and moreover, these are \textit{exact.} The category $\mathrm{Ho}_W(\M) \simeq \mathrm{Ho}_W(\Mc)$ ``is computed'' by a calculus of left fractions in $\pi\Mc$, where $\pi\Mc$ is the category obtained from $\Mc$ by replacing the $\mathrm{Hom}_{\Mc}(X,Y) = \mathrm{Hom}_{\M}(X,Y)$ in it by the quotient by a suitable homotopy relation:
\begin{description}[leftmargin=9em, labelwidth=8em, font=\normalfont]
\item[$f,g : X \rightrightarrows Y$] are QB-homotopic (in the sense of \textit{Quillen-Brown}) if there exists an $X' \xrightarrow{\alpha} X$ in $W$ (and we may assume $X' \in \mathrm{Ob}\,\Mc$ as well) such that $f\alpha \overset{1}{\sim} g\alpha$ in the sense of \textsc{Quillen}, i.e. there exists a cylinder object
\begin{tikzcd}[baseline=(current bounding box.center)]
I_{X'} \arrow[r, "\in W"] & X' \\
X' \vee X' \arrow[u, "\in C"] \arrow[ur, dotted] &
\end{tikzcd}
for $X'$, such that $(f\alpha, g\alpha)\ :\ X' \vee X' \longrightarrow Y$ factors through an $h' : I_{X'} \longrightarrow Y$.
\end{description}

It is from this criterion that the (left) exactness of $W$-cocartesian squares follows.

I want to see how a cofibration structure on $\M$ extends to the categories $\M(I^\circ) = \underline{\mathrm{Hom}}(I, M)$. We have in $\M(I^\circ)$ \msp{42} the set $W_{I^\circ}$, and $C_I$ (it is probably better to denote it also by $C_{I^\circ}$) defined in the preceding paragraph. It is a matter of checking the axioms. The axiom $\mathrm{Cof}_0$ is immediate, and well known for $I$ an arbitrary category. % [sic] \underline{Hom}(I, M) -> \underline{Hom}(I, \mathcal{M}), italic M printed where the script \mathcal{M} is used everywhere else

$\mathrm{Cof}_1$ : $C_{I^\circ}$ contains the isomorphisms, that is clear, because $C$ contains them. Stable under composition $X_\bullet \xrightarrow{u} Y_\bullet \xrightarrow{v} Z_\bullet$, let us look at $K \subset J$, put $w = v \circ u$,
\[
\begin{tikzcd}[execute at end picture={\draw[decorate,decoration={brace,mirror,amplitude=6pt}] ([yshift=-2pt]B.south west) -- ([yshift=-2pt]E.south east) node[midway,below=7pt]{$\scriptstyle w_{K,J}$};}]
X_K \arrow[r] \arrow[d] & X_J \arrow[d] & & \\
Y_K \arrow[r] \arrow[d] & Y_{K,J} \arrow[r, "u_{K,J}"] \arrow[d] & Y_J \arrow[d] & \\
Z_K \arrow[r] & |[alias=B]| \bullet \arrow[r, "u'_{K,J}"] & Z_{K,J} \arrow[r, "v_{K,J}"] & |[alias=E]| Z_J \rlap{\ ,}
\end{tikzcd}
\]
\vspace{1ex}

$w_{K,J}$ is the composite of $v_{K,J}$, which is in $C$, and $u'_{K,J}$, which is too since obtained from $u_{K,J}$ by cobase change.

(\textbf{N.B.} We have already made the argument with a single $u_\bullet$, and \textit{three} open sets $J \subset K \subset L$. It \msp{43} would be appropriate to introduce the notion of a $C$-square
\begin{tikzcd}[baseline=(current bounding box.center)]
X_0 \arrow[r] \arrow[d] & X_1 \arrow[d] \\
X_2 \arrow[r] & X_3
\end{tikzcd}, as a square such that
\printedpage{23}
% head: A. Grothendieck : Dérivateurs, Ch. XVII, Édition des Universités de Montpellier II et Paris VII
$X_1 \amalg_{X_0} X_2 \longrightarrow X_3$ is in $C$ (\textbf{N.B.} notion symmetric for the symmetry functors of squares), and say that the composite of two $C$-squares is a $C$-square.)

Cobase change
\[
\begin{tikzcd}
X_\bullet \arrow[r, "u"] \arrow[d] & Y_\bullet \arrow[d] \\
X'_\bullet \arrow[r, "u'"] & Y'_\bullet \rlap{\ .}
\end{tikzcd}
\]

Let us consider $K \subset J$, we have
\[
\begin{tikzcd}[row sep=large, column sep=large]
& X_K \arrow[r] \arrow[d] \arrow[ddl] \arrow[dr, phantom, "\text{\scriptsize cocart.}"] & X_J \arrow[d] \arrow[dr] & \\
& Y_K \arrow[r] \arrow[ddl] \arrow[dl, phantom, "\text{\scriptsize cocart.}"] & Y_{K,J} \arrow[r, "u_{K,J}"] \arrow[ddl] \arrow[d, phantom, "\text{\scriptsize cocartesian}"] & Y_J \arrow[ddl] \\
X'_K \arrow[r, crossing over] \arrow[d] \arrow[dr, phantom, "\text{\scriptsize cocart.}"] & X'_J \arrow[d] \arrow[dr, crossing over] \arrow[from=1-3, crossing over] & {} & \\
Y'_K \arrow[r] & Y'_{K,J} \arrow[r, "u'_{K,J}"] & Y'_J \rlap{\ ,} &
\end{tikzcd}
\]
where the three lateral squares marked ``cocart.'' are cocartesian, as is the square in an inclined plane. It follows, first that the fourth lateral square of the cube is itself also cocartesian, \msp{44} then that the square having two sides $u_{K,J}$, $u'_{K,J}$ is cocartesian. (\textbf{N.B.} compare the similar diagram p. 32 -- there would be a common lemma to extract, distinguishing two cases \ldots)

$\mathrm{Cof}_2$. To prove that $C_{I^\circ} \cap W_{I^\circ}$ is stable under cobase change, I see no way of doing it, if I do not already know that we have
\[ C_{I^\circ} \cap W_{I^\circ} \ = \ (TC)_{I^\circ}\ , \tag*{$(*)$} \]
and then we apply the preceding result (where $W$ did not enter) to $TC$ instead of $C$. But to have $(*)$, I must assume that we are in one of the following two cases (so as to be able to apply theorem 2.1 and its corollary 2.5).
\begin{enumerate}[leftmargin=2.5em]
\item[a)] $I$ finite, or more generally $I$ artinian \textit{and} $\mathrm{Top}(I)$ is noetherian.
\item[b)] $I$ artinian, and moreover $C$ and $TC$ are stable under transfinite composition.
\end{enumerate}
$\mathrm{Cof}_3$. The factorization axiom seems to me a priori the most delicate\footnote[15]{The argument that follows applies to other factorization situations, for example for a given Quillen precouple $(\Phi, \Psi)$ in $\M$.}. But in the end this comes out without too much trouble. Let there be given % [sic] facorisation -> factorisation, missing letter
\[ f \ : \ X_\bullet \longrightarrow Y_\bullet \]

\end{document}
